\section{Case Study for ToG-3}
\label{appendix:case_study}

This section provides a detailed case study of ToG-3 in deep reasoning task (Figure~\ref{fig:case_study_evolving_subgraph}) and broad reasoning task (Figure~\ref{fig:case_study_regression_metrics} and Figure~\ref{fig:case_study_food_policy}), offering an intuitive demonstration of the execution dynamics of its dual-evolution mechanism—comprising Evolving Query and Evolving Subgraph—across multi-step reasoning processes.


\begin{figure*}
\centering
\begin{tikzpicture}[font=\sffamily, scale=0.88, every node/.style={scale=0.88}]
\node[anchor=north] at (0,0) {
\begin{tcolorbox}[
    width=16cm,
    colback=gray!5,
    colframe=gray!40,
    arc=3mm,
    boxrule=1pt,
    left=8pt,
    right=8pt,
    top=6pt,
    bottom=6pt,
    fontupper=\ttfamily\small,
    breakable
]
{\bfseries\color{roleblue} Question:} \\
{\color{varcolor}What nationality is the performer of the song When The Stars Go Blue?}

\vspace{8pt}
{\bfseries\color{roleblue} Initial Evidence (Sub-Graph):} \\
{\color{evicolor}When The Stars Go Blue -> performed\_by -> Ryan Adams} \\
{\color{evicolor}Ryan Adams -> known\_for -> alternative country, rock, folk} \\
{\color{evicolor}Ryan Adams -> also\_includes\_genre -> indie rock, Americana} \\
{\color{evicolor}... (no nationality information present)}

\vspace{8pt}
{\bfseries\color{roleblue} Evolving Query Decomposition 1 (Reflector Agent)} \\
{\color{varcolor} Who is the performer of the song When The Stars Go Blue?
}\\
{\color{anscolor} Ryan Adams.
}

\vspace{8pt}
{\bfseries\color{roleblue} Evolving Sub-Graph Refinement (Constructor Agent):} \\
{\color{evicolor}Sammy Hagar -> is -> American rock vocalist} \\
{\color{evicolor}Chuck Berry -> is -> American singer and songwriter} \\
{\color{evicolor}Elvis Presley -> pioneered -> Rock and Roll (American origin)} \\
{\color{evicolor}\textbf{Ryan Adams -> nationality -> American} \quad [Added]} \\
{\color{evicolor}Wee Willie Walker -> was -> American soul and blues singer}

\vspace{8pt}
{\bfseries\color{roleblue} Evolving Query Decomposition 2 (Reflector Agent)} \\
{\color{varcolor} What is Ryan Adams known for besides the song When The Stars Go Blue?
}\\
{\color{anscolor}  Ryan Adams is known for his work as a singer-songwriter and musician, with a musical style that encompasses alternative country, rock, and folk.
}

\vspace{8pt}
{\bfseries\color{roleblue} Evolving Query Decomposition 3 (Reflector Agent)} \\
{\color{varcolor} What is Ryan Adams's nationality?
}\\
{\color{anscolor}  Ryan Adams is an American singer and musician.
}

\vspace{8pt}
{\bfseries\color{roleblue} Final Answer (Response Agent):} \\
{\color{anscolor}ryan adams is american.}

\vspace{8pt}
{\bfseries\color{roleblue} Ground Truth:} \\
{\color{evalcolor}america}

\vspace{6pt}
{\bfseries\color{roleblue} Explanation:} \\
{\color{evalcolor}
The initial sub-graph contained no explicit information about Ryan Adams' nationality.\\
Through evolving query decomposition, the system identified a missing link and triggered sub-graph refinement.\\
New evidence from biographical triples (e.g., "Ryan Adams -> nationality -> American") was integrated, enabling a definitive answer.
}
\end{tcolorbox}
};
\end{tikzpicture}
\caption{Case Study: Evolving Query and Evolving Sub-Graph for Multi-Hop Reasoning. This example demonstrates how an initial knowledge sub-graph fails to answer a complex question, prompting a decomposition into a sub-query. The Constructor Agent refines the sub-graph with additional biographical facts, enabling the Response Agent to generate the correct answer. The process highlights the dynamic, iterative nature of self-evolving context retrieval.}
\label{fig:case_study_evolving_subgraph}
\end{figure*}

\begin{figure*}
\centering
\begin{tikzpicture}[font=\sffamily, scale=0.88, every node/.style={scale=0.88}]
\node[anchor=north] at (0,0) {
\begin{tcolorbox}[
    width=16cm,
    colback=gray!5,
    colframe=gray!40,
    arc=3mm,
    boxrule=1pt,
    left=8pt,
    right=8pt,
    top=6pt,
    bottom=6pt,
    fontupper=\ttfamily\small,
    breakable
]
{\bfseries\color{roleblue} Question:} \\
{\color{varcolor}How do different regression metrics (e.g., RMSE, MAE, R\textsuperscript{2}) vary by model and dataset?}

\vspace{8pt}
{\bfseries\color{roleblue} Evidence (from ToG-3):} \\
{\color{evicolor}Mean Squared Error -> is\_related\_to -> Root Mean Squared Error...} \\
{\color{evicolor}Mean Squared Error -> is\_similar\_to -> Mean Absolute Error...} \\
{\color{evicolor}summary\_23 -> SUMMARY\_FOR -> RMSE.MSE emphasizes larger errors...} \\
{\color{evicolor}summary\_287 -> SUMMARY\_FOR -> MAE. ARIMA model performance evaluated using MAE...} \\

\vspace{8pt}
{\bfseries\color{roleblue} Answer 1(from ToG-3):} \\
{\color{anscolor}Different regression metrics such as RMSE, MAE, and R\textsuperscript{2} can vary by model and dataset...} \\
{\color{anscolor}RMSE penalizes large errors more heavily due to squared terms, while MAE is robust to outliers...} \\
{\color{anscolor}R\textsuperscript{2} measures explained variance but can be misleading on small or overfitted datasets...} \\

\vspace{8pt}
{\bfseries\color{roleblue} Answer 2(from GraphRAG):} \\
{\color{anscolor}The provided community summary does not detail specific variations...To evaluate how these metrics vary:} \\
{\color{anscolor}1. Train different regression models on your datasets...} \\
{\color{anscolor}2. Compute RMSE, MAE, R\textsuperscript{2} using MLlib.RegressionMetrics...} \\
{\color{anscolor}3. Compare performance across models and datasets...} \\


\vspace{6pt}
{\bfseries\color{roleblue} Evaluation Result (JSON):}
{\bfseries\color{evalcolor}
\begin{verbatim}
{
  "Comprehensiveness": {
    "Winner": "Answer 1",
    "Explanation": "Answer 1 provides a detailed explanation of various regression metrics..."
  },
  "Diversity": {
    "Winner": "Answer 2",
    "Explanation": "Answer 2 presents a different perspective by incorporating specific tools..."
  },
  "Empowerment": {
    "Winner": "Answer 1",
    "Explanation": "Answer 1 empowers the reader by explaining what each metric means..."
  },
  "Overall Winner": {
    "Winner": "Answer 1",
    "Explanation": "Answer 1 is the overall winner as it provides a comprehensive understanding..."
  }
}
\end{verbatim}
}
\end{tcolorbox}
};
\end{tikzpicture}
\caption{Case Study: Comparing Regression Metrics Across Models and Datasets. This example illustrates how two reasoning systems answer a technical ML question: GraphRAG emphasizes practical implementation (e.g., using Spark's MLlib), while ToG3 focuses on theoretical distinctions between RMSE, MAE, and R\textsuperscript{2}. An evaluator selects the more comprehensive and empowering answer based on evidence from the knowledge graph.}
\label{fig:case_study_regression_metrics}
\end{figure*}



\begin{figure*}
\centering
\begin{tikzpicture}[font=\sffamily, scale=0.88, every node/.style={scale=0.88}]
\node[anchor=north] at (0,0) {
\begin{tcolorbox}[
    width=16cm,
    colback=gray!5,
    colframe=gray!40,
    arc=3mm,
    boxrule=1pt,
    left=8pt,
    right=8pt,
    top=6pt,
    bottom=6pt,
    fontupper=\ttfamily\small,
    breakable
]
{\bfseries\color{roleblue} Question:} \\
{\color{varcolor}What policy changes are needed to foster equitable access to nutritious food?}

\vspace{8pt}
{\bfseries\color{roleblue} Evidence (from ToG-3):} \\
{\color{evicolor}summary\_216 -> SUMMARY\_FOR -> Food Security. **Supermarket and Food Security**: The reliance on supermarkets for food is viewed as a threat...} \\
{\color{evicolor}Victory Garden -> promoted -> Food Security... Supermarket -> threatens -> Food Security..."} \\
{\color{evicolor}White House Childhood Obesity Task Force -> recommended -> Farm to School...} \\
{\color{evicolor}Food Works -> collaborates\_with -> Sauvie Island Organics...} \\
{\color{evicolor}Michael Taylor -> interprets -> Food Modernization Act... background in representing Monsanto...} \\

{\bfseries\color{roleblue} Answer 1(from ToG-3):} \\
{\color{anscolor}policy changes that could foster equitable access to nutritious food may include implementing subsidies...} \\
{\color{anscolor}increasing funding for farm to school programs in underserved communities...} \\
{\color{anscolor}establishing regulations that prioritize the availability of nutritious options in school meal programs...} \\
{\color{anscolor}and creating educational initiatives to promote awareness about nutrition and healthy eating habits...} \\
{\color{anscolor}these changes could help address disparities in access to nutritious food...}\\

\vspace{8pt}
{\bfseries\color{roleblue} Answer 2(from GraphRAG):} \\
{\color{anscolor}1. **strengthen conflict of interest policies**: prevent individuals with significant industry ties...} \\
{\color{anscolor}2. **enhance transparency and accountability**: implement transparent processes...} \\
{\color{anscolor}3. **support diverse food producers**: promote a variety of food producers...} \\
{\color{anscolor}4. **expand nutrition education and assistance programs**: enhance nutrition education...} \\
{\color{anscolor}5. **ensure regulatory independence**: guarantee that regulatory bodies are independent...} \\
{\color{anscolor}6. **promote local food systems**: encourage local food production through grants, subsidies...} \\
{\color{anscolor}7. **support community gardens and urban agriculture**: provide funding and resources...} \\
{\color{anscolor}8. **promote sustainable farming practices**: offer incentives for sustainable farming...} \\
{\color{anscolor}These measures aim to create a more resilient and sustainable food system...}


\vspace{6pt}
{\bfseries\color{roleblue} Evaluation Result (JSON):}
{\bfseries\color{evalcolor}
\begin{verbatim}
{
{
  "Comprehensiveness": {
    "Winner": "Answer 2",
    "Explanation": "Answer 2 provides a more detailed and structured set of policy changes... "
    "broader scope than Answer 1..."
  },
  "Diversity": {
    "Winner": "Answer 2",
    "Explanation": "Answer 2 presents a wider variety of perspectives... "
    "including regulatory, community, and economic solutions."
  },
  "Empowerment": {
    "Winner": "Answer 1",
    "Explanation": "Answer 1 empowers readers by offering a comprehensive overview..." 
    "facilitates informed judgment."
  },
  "Overall Winner": {
    "Winner": "Answer 2",
    "Explanation": "Answer 2 emerges as the overall winner due to its superior"
    "comprehensiveness, diversity, and empowerment."
  }
}
\end{verbatim}
}
\end{tcolorbox}
};
\end{tikzpicture}
\caption{Case Study: Policy Recommendations for Equitable Food Access. This example illustrates the full reasoning pipeline: a complex policy question is answered by two different systems (GraphRAG and ToG-3), supported by retrieved knowledge snippets. An evaluator then compares both responses across multiple dimensions, selecting the more comprehensive, diverse, and empowering answer as the winner.}
\label{fig:case_study_food_policy}
\end{figure*}